\documentclass[11pt]{article}
\usepackage[margin=1in]{geometry}
\usepackage{amsmath,amssymb,mathtools}
\usepackage{fontspec,unicode-math}
\setmainfont{Latin Modern Roman}
\setsansfont{Latin Modern Sans}
\setmathfont{Latin Modern Math}
\usepackage{xurl,hyperref}
\hypersetup{colorlinks=true,urlcolor=blue,linkcolor=blue}
\newcommand{\sref}[2]{\hyperlink{source:#1:#2}{[S#2]}}
\newcommand{\cataloguescope}[1]{\paragraph{Recorded mathematical question.}#1}
\begin{document}
% UnsolvedMath ID 10000066; source catalogue code AMR-099-0066
\setcounter{section}{10000065}
\section{Fluctuations in recursive hierarchical first-passage percolation}\label{10000066}
{\small\sffamily UnsolvedMath ID 10000066 · Probability}
\subsection{Definitions and mathematical statement}

\paragraph{Shared notation.}$\mathbb N=\{1,2,\ldots\}$, and $\mathbb Z,\mathbb Q,\mathbb R,\mathbb C$ denote the integers, rationals, reals and complex numbers. Any other conventions are specified locally.


Let $G_0$ be one oriented edge from a source to a sink. Replace each edge recursively by a copy of the three-edge pattern with a single edge from its source to an intermediate junction and two parallel edges from that junction to its sink, preserving the source-to-sink orientation; call the result $G_n$. Assign independent positive weights with common law $\xi$ to its edges and let $D_n$ be the minimum total weight of a source--sink path. Thus $D_0\stackrel d=\xi$ and
\[D_n\stackrel d=D_{n-1}^{(1)}+\min\{D_{n-1}^{(2)},D_{n-1}^{(3)}\},\]
where the three variables on the right are independent copies. For the comparison branching random walk, a binary tree starting with an edge has independent edge weights of the same law; let $M_m$ be its minimum root-to-level-$m$ weight. For laws for which the indicated finite limits exist, define the constants
\[\gamma_{\rm rec}=\lim_{n\to\infty}2^{-n}\mathbb E D_n,\qquad\gamma_{\rm brw}=\lim_{m\to\infty}m^{-1}\mathbb E M_m.\]
\paragraph{Statement recorded in the supplied source.}
Determine the scale and limiting behavior of $D_n-\mathbb E D_n$ and obtain lower-order terms beyond the leading growth $\mathbb E D_n\sim 2^n\gamma_{\rm rec}$ when this finite constant exists. Determine the distributional and integrability hypotheses on $\xi$ under which the requested asymptotics hold. In particular, identify the concentration and lower-order corrections corresponding to those for the minimum position of the comparison branching random walk. \sref{10000066}{2}
\subsection{Short English statement}
Find the fluctuations and lower-order growth terms of hierarchical first-passage distances.
\subsection{Sources}
\begin{description}
\item[\hypertarget{source:10000066:1}{S1}] ULAM.AI and the UnsolvedMath Contributors, \emph{UnsolvedMath: A Curated Collection of Open Mathematics Problems}, record 10000066 (AMR-099-0066). CC BY 4.0. \url{https://huggingface.co/datasets/ulamai/UnsolvedMath}.
\item[\hypertarget{source:10000066:2}{S2}] Itai Benjamini, \emph{Euclidean vs. Graph Metric}, version dated 16 July 2012; original author manuscript preserved in Arquivo.pt capture of 31 December 2020, Section 9.2, Question 9.2, p. 15. \url{https://arquivo.pt/noFrame/replay/20201231041538id_/http://www.wisdom.weizmann.ac.il/~itai/erd100.pdf}.
\end{description}
\label{end:10000066}
\end{document}
